% TOP_PROBLEM: {"id": 20000913, "title": "Random sub-neighborhood discrepancy on the Boolean cube", "category_id": 2, "category": {"id": 2, "name": "combinatorics", "display_name": "Combinatorics", "description": "Counting problems, graph theory, discrete structures.", "slug": "combinatorics", "order_index": 2, "created_at": "2026-07-31T15:26:25.671Z"}, "status": "partially_solved", "problem_number": "AIM-COMBINATORICS-0038", "set_id": 14, "statusQualification": "Random subsets are interpreted as independent uniform subsets, one per cube vertex."}
% FIRST_POSTED: 2026-10-01
% ENTRY_KIND: statement_only
% UnsolvedMath ID 20000913; source catalogue code AIM-COMBINATORICS-0038
% !TeX program = lualatex
% Definition and sources only; this file is not an LLM solution attempt.
\documentclass[11pt,a4paper]{article}
\usepackage[margin=25mm]{geometry}
\usepackage{fontspec}
\setmainfont{lmroman10-regular.otf}[BoldFont=lmroman10-bold.otf,
 ItalicFont=lmroman10-italic.otf,BoldItalicFont=lmroman10-bolditalic.otf]
\usepackage{amsmath,amssymb,mathtools}
\usepackage{unicode-math}
\setmathfont{latinmodern-math.otf}
\AtBeginDocument{\let\setminus\smallsetminus}
\usepackage{xurl,hyperref}
\hypersetup{colorlinks=true,urlcolor=blue}
\setcounter{secnumdepth}{0}
\setlength{\parindent}{0pt}
\setlength{\parskip}{5pt}
\setlength{\emergencystretch}{3em}
\begin{document}
\section*{Random sub-neighborhood discrepancy on the Boolean cube}\label{20000913}
{\small UnsolvedMath ID 20000913 · AIM-COMBINATORICS-0038 · Combinatorics · AIM Workshop Problem Lists}
\subsection{Definitions and mathematical statement}
For $n\ge1$, let $Q_n$ be the graph with vertex set $V_n=\{0,1\}^n$,
where two vertices are adjacent when they differ in exactly one
coordinate. For each $v\in V_n$, choose a subset $S_v$ of its $n$
neighbours uniformly from all $2^n$ subsets. Make these choices
independently for different vertices. Equivalently, each ordered
neighbour pair $(v,u)$ independently has probability $1/2$ of placing
$u$ in $S_v$. In particular, membership of $u$ in $S_v$ need not agree
with membership of $v$ in $S_u$.

For the resulting random set system $\mathcal S_n=(S_v:v\in V_n)$,
its discrepancy is the random variable
\[
 D_n=\min_{\varepsilon:V_n\to\{-1,1\}}
        \max_{v\in V_n}\left|\sum_{u\in S_v}\varepsilon(u)\right|.
\]
The signing is chosen after the entire set system has been sampled and
is shared by all its sets. Empty sets contribute zero to the maximum.
The problem is to determine the typical asymptotic size of $D_n$ as
$n\to\infty$, ideally by matching upper and lower bounds holding with
probability tending to one. For example, an order-of-magnitude answer
would specify a deterministic function $f(n)$ and constants $c,C>0$
for which $\Pr(c f(n)\le D_n\le C f(n))\to1$.
The independence and uniform-subset conventions make the workshop's
phrase ``random subset'' precise.
\subsection{Short English statement}
How well can one balance independently chosen random subsets of the neighbourhoods of a high-dimensional cube?
\subsection{Sources}
\begin{itemize}
\item[S1] ULAM.AI and the UnsolvedMath Contributors, supplied problems.json, record 20000913 (AIM-COMBINATORICS-0038), AIM Workshop Problem Lists. Catalogue identity and source transcription; CC BY 4.0.
\url{https://huggingface.co/datasets/ulamai/UnsolvedMath}.
\item[S2] American Institute of Mathematics, Hereditary discrepancy and factorization norms, item 1.25. Original workshop question underlying AIM-COMBINATORICS-0038.
\url{http://aimpl.org/hereddiscrep/1/}.
\end{itemize}
\end{document}
